\section{Introduction}

Large language models are increasingly deployed in settings where they must
choose between disagreeing knowledge sources. A model may have a strong
parametric prior, receive retrieved evidence that contradicts that prior, and
then spend additional reasoning tokens reinforcing one side of the conflict.
Current systems still treat this arbitration problem mostly heuristically.
Adaptive RAG systems decide when to retrieve or how much to trust retrieval
using tuned thresholds or learned controllers, while conflict benchmarks
typically evaluate behavior descriptively rather than against a normative
decision rule \citep{xu2024knowledgeconflicts,mallen2023popqa,asai2024selfrag}.
At the same time, recent calibration papers have split into two apparently
contradictory camps: on standard QA without explicit knowledge conflict,
reasoning models often become \emph{better} calibrated as CoT unfolds
\citep{yoon2025confidence}, while in harder or misspecified settings longer
reasoning can instead drive systematic overconfidence
\citep{lacombe2025dontthinktwice,welch2026cost,damani2025rlcr}.

This paper takes a different view. We treat knowledge conflict as a
\emph{posterior-predictive decision problem}. The core object is not ``should
the model use retrieval,'' but rather ``what predictive distribution minimizes
Bayes risk given a parametric prior, a context-conditioned likelihood, and a
latent reliability variable for the context?'' Once written this way, two
theorem-level questions and one structurally motivated reasoning-time question
become precise:

\begin{enumerate}
\item What is the Bayes-optimal arbitration rule?
\item How bad are policies that always trust one source or use one global mix
weight?
\item How does longer reasoning interact with calibration when the two sources
conflict?
\end{enumerate}

These questions are motivated by a real gap in the literature. Surveys and
benchmarks on knowledge conflict repeatedly argue that blindly prioritizing
either parametric memory or context is undesirable, but they stop short of a
formal regret or minimax statement \citep{xu2024knowledgeconflicts}. Empirical
systems such as PopQA-style adaptive retrieval, Self-RAG, FLARE, and related
adaptive-RAG methods use thresholds or controllers that are not derived from a
Bayes optimum \citep{mallen2023popqa,asai2024selfrag}. Meanwhile, the most
relevant calibration literature shows that longer reasoning or more test-time
compute can worsen calibration, but the strongest claims either live outside
knowledge conflict or remain descriptive rather than decision-theoretic
\citep{halawi2024overthinking,chen2025rethinking,ghosal2025mirage,
sun2025context,yoon2025confidence,lacombe2025dontthinktwice,welch2026cost,
damani2025rlcr}.

Our contributions are threefold, but not all at the same level of formal
rigor. First, we derive a Bayes-style
reliability-aware arbitration rule and validate it on a corrected broad wave of
benchmark-backed experiments. In that wave, the Bayes proxy achieves mean
regret $-0.0461$ versus $-0.0233$ for a generic adaptive heuristic, while
`always\_context' and `always\_parametric' degrade to $7.2237$ and $5.9356$.
Second, we show that fixed trust policies are minimax-bad in practice on
conflict-heavy slices: in the conflict wave, the Bayes proxy reaches
$-0.1256$ versus $-0.0752$ for the heuristic, while fixed trust policies remain
at $5.9037$ and $7.1329$. Third, we revisit the reasoning-time calibration
claim using real DeepSeek-R1-Distill traces rather than proxy CoT buckets.
This third contribution is best read as a structurally motivated
misspecification account, grounded in generalized-Bayes and Berk--Nash
analogies plus direct trace diagnostics, rather than as a closed-form theorem
of the same status as the first two results. The result is stronger and more
nuanced than our original monotone hypothesis:
reasoning amplifies overconfidence under conflict, but the shape depends on the
conflict family. On WikiContradict, both 7B and 14B show an intermediate-CoT
peak followed by partial recovery. On ConflictBank conflict, 7B shows the same
peak-and-recover shape, while 14B remains severely overconfident through long
CoT. A completed same-family Qwen sweep sharpens the interpretation further:
the specific ``7B recovers / 14B saturates'' asymmetry does not replicate
universally across families, but a benchmark-dependent two-regime law does.
Qwen shows no recovery through 32B on controlled ConflictBank conflict, while
32B is the first scale at which recovery reappears on naturalistic
WikiContradict contradiction. The cleanest interpretation is not a universal
scale law but a conditional misspecification law: some reasoning-trained models
and conflict families enter a regime where longer CoT sharpens confidence
faster than it improves accuracy.

The paper therefore makes two negative claims and one positive one. The
negative claims are that fixed trust in retrieval is wrong, and fixed trust in
the parametric model is wrong. The positive claim is that knowledge arbitration
admits a principled reliability-aware rule that improves regret, and that
reasoning-time confidence under conflict follows a structured, measurable law
rather than arbitrary noise.

We also make three scope boundaries explicit up front. First, our current
training-time comparison to RLCR-style calibration improvement is not yet a
head-to-head number in the paper body; the present intervention evidence is
confidence-only $\eta$-tempering plus a smaller training-time pilot, not a full
Damani-style reward redesign. Second, the retrieval deployment evidence is a
bounded Wikipedia-backed demo rather than a production-scale retrieval stack.
Third, the current experiments are English-only, so the reported gains should
not be read as multilingual calibration guarantees.
